\clearpage
\thispagestyle{plain}
\markboth{PRELIMINARIES}{PRELIMINARIES}
\oldpage[IV-1]{5}
\phantomsection
\label{chapter:0IV1}

\begin{center}
{\large CHAPTER 0 (continued)\footnotemark\par}
\vspace{1.5em}
{\Large\bfseries PRELIMINARIES\par}
\vspace{2em}
{\large\bfseries Summary\par}
\end{center}

\footnotetext{Recall that the symbol $0_N$ (resp. $\operatorname{Err}_N$)
refers to the part of Chapter 0 (resp. the list of Errata) attached to Chapter
$N$.}

\begin{center}
\begin{tabular}{@{}r l@{}}
\S~14. & Combinatorial dimension of a topological space.\\
\S~15. & M-regular sequences and $\mathscr F$-regular sequences.\\
\S~16. & Dimension and depth in Noetherian local rings.\\
\S~17. & Regular rings.\\
\S~18. & Supplement on extensions of algebras.\\
\S~19. & Formally smooth algebras and Cohen rings.\\
\S~20. & Derivations and differentials.\\
\S~21. & Differentials in rings of characteristic $p$.\\
\S~22. & Differential criteria for formal smoothness and regularity.\\
\S~23. & Japanese rings.
\end{tabular}
\end{center}

Almost all of the preceding sections are devoted to the exposition of
notions of Commutative Algebra that will be used in the course of Chapter IV.
Although a good part of these notions already appears in several works
([1], [12], [13], [17], [30]), we thought that it would be more convenient for the
reader to have a continuous and more or less self-contained exposition of them.
Together with \S\S~5, 6 and 7 of Chapter IV (where the language of schemes is
used), these sections constitute, within our Treatise, a small special
Treatise, more or less independent of Chapters I to III, which aims to present
in a coherent form the properties of rings that ``behave well''
with respect to operations such as completion or integral
closure, by systematically relating these properties to a few
general concepts\footnote{Most of the properties in question are due to Chevalley, Zariski, Nagata and Serre. The method followed here was first set out by A.~Grothendieck in a course at Harvard University in the autumn of 1961.}.

\setcounter{section}{13}
\section{Combinatorial dimension of a topological space}
\label{section:0.14}

\setcounter{subsection}{0}
\subsection{Combinatorial dimension of a topological space}
\label{subsection:0.14.1}

\oldpage[0\textsubscript{IV-1}]{6}
\begin{env}[14.1.1]
\label{0.14.1.1}
Let $I$ be an ordered set; a \emph{chain} of elements of $I$ is by
definition a strictly increasing finite sequence
$i_0<i_1<\cdots<i_n$ of elements of $I$ ($n\geq 0$); by definition the
\emph{length} of this chain is $n$. If $X$ is a topological space,
the set of its \emph{irreducible closed} subsets is ordered by
inclusion, whence the notion of a \emph{chain} of irreducible closed subsets of
$X$.
\end{env}

\begin{definition}[14.1.2]
\label{0.14.1.2}
Let $X$ be a topological space. We call \emph{combinatorial dimension of
$X$} (or simply \emph{dimension of $X$} if no confusion can
arise), and denote by $\dim_c(X)$, or simply $\dim(X)$, the
supremum of the lengths of the chains of irreducible closed subsets of $X$.
For every $x\in X$, we call \emph{combinatorial dimension of $X$ at $x$}
(or simply \emph{dimension of $X$ at $x$}), and denote by $\dim_x(X)$, the
number $\inf_U(\dim(U))$, where $U$ runs through the set of open neighbourhoods of
$x$ in $X$.
\end{definition}

It follows from this definition that we have
\[
  \dim(\varnothing)=-\infty
\]
(the supremum in $\overline{\mathbb R}$ of the empty subset being
$-\infty$). If $(X_\alpha)$ is the family of irreducible components of
$X$, we have
\[
  \dim(X)=\sup_\alpha\dim(X_\alpha)
  \tag{14.1.2.1}\label{0.14.1.2.1}
\]
since every chain of irreducible closed subsets of $X$ is by definition
contained in an irreducible component of $X$, and conversely these
components are closed in $X$, so every irreducible closed subset of an
$X_\alpha$ is an irreducible closed subset of $X$.

\begin{definition}[14.1.3]
\label{0.14.1.3}
We say that a topological space $X$ is \emph{equidimensional} if all its
irreducible components have the same dimension (then equal to $\dim(X)$ by
(14.1.2.1)).
\end{definition}

\begin{proposition}[14.1.4]
\label{0.14.1.4}
\begin{enumerate}
  \item[\rm(i)] For every subset $Y$ of a topological space $X$, we have
  $\dim(Y)\leq\dim(X)$.
  \item[\rm(ii)] If a topological space $X$ is a finite union of closed
  subsets $X_i$, we have $\dim(X)=\sup_i\dim(X_i)$.
\end{enumerate}
\end{proposition}

\begin{proof}
For every irreducible closed subset $Z$ of $Y$, the closure $\overline Z$ of
$Z$ in $X$ is irreducible ($0_{\mathrm I}$, 2.1.2) and
$\overline Z\cap Y=Z$, whence (i). On the other hand, if
$X=\bigcup_{i=1}^n X_i$, where the $X_i$ are closed, every irreducible closed
subset of $X$ is contained in one of the $X_i$ ($0_{\mathrm I}$, 2.1.1),
and consequently every chain of irreducible closed subsets of $X$ is contained
in one of the $X_i$; whence (ii).
\end{proof}

We deduce from (14.1.4, (i)) that for every $x\in X$ we can also write
\[
  \dim_x(X)=\lim_U\dim(U)
  \tag{14.1.4.1}\label{0.14.1.4.1}
\]
the limit being taken along the decreasing filtered set of open
neighbourhoods of $x$ in $X$.

\oldpage[0\textsubscript{IV-1}]{7}
\begin{corollary}[14.1.5]
\label{0.14.1.5}
Let $X$ be a topological space, $x$ a point of $X$, $U$ a neighbourhood of
$x$, and $Y_i$ ($1\leq i\leq n$) closed subsets of $U$ such that
$x\in Y_i$ for every $i$ and $U$ is the union of the $Y_i$. Then we have
\[
  \dim_x(X)=\sup_i\bigl(\dim_x(Y_i)\bigr).
  \tag{14.1.5.1}\label{0.14.1.5.1}
\]
\end{corollary}

\begin{proof}
It follows from (14.1.4, (ii)) that we have
$\dim_x(X)=\inf_V(\sup_i(\dim(Y_i\cap V)))$ when $V$ runs through the set
of open neighbourhoods of $x$ contained in $U$; likewise we have
$\dim_x(Y_i)=\inf_V(\dim(Y_i\cap V))$ for every $i$. The corollary is therefore
evident if
\[
  \sup_i\bigl(\dim_x(Y_i)\bigr)=+\infty~;
\]
otherwise, there is an open neighbourhood $V_0\subset U$ of $x$ such
that $\dim(Y_i\cap V)=\dim_x(Y_i)$ for $1\leq i\leq n$ and for every
$V\subset V_0$, whence the conclusion.
\end{proof}

\begin{proposition}[14.1.6]
\label{0.14.1.6}
For every topological space $X$, we have
$\dim(X)=\sup_{x\in X}\dim_x(X)$.
\end{proposition}

\begin{proof}
It follows from the definition (14.1.2) and from Proposition (14.1.4) that
$\dim_x(X)\leq\dim(X)$ for every $x\in X$. On the other hand, let
$Z_0\subset Z_1\subset\cdots\subset Z_n$ be a chain of irreducible closed
subsets of $X$, and let $x\in Z_0$; for every open subset $U\subset X$
containing $x$, $U\cap Z_i$ is irreducible ($0_{\mathrm I}$, 2.1.6) and
closed in $U$, and since we have $\overline{U\cap Z_i}=Z_i$ in $X$, the
$U\cap Z_i$ are pairwise distinct; therefore we have $\dim(U)\geq n$, which
completes the proof.
\end{proof}

\begin{corollary}[14.1.7]
\label{0.14.1.7}
If $(X_\alpha)$ is an open covering of $X$, or a locally finite closed
covering of $X$, we have
$\dim(X)=\sup_\alpha(\dim(X_\alpha))$.
\end{corollary}

\begin{proof}
If $X_\alpha$ is a neighbourhood of $x\in X$, we have
$\dim_x(X)\leq\dim(X_\alpha)$, whence the first assertion. On the other hand, if
the $X_\alpha$ are closed and if $U$ is a neighbourhood of $x\in X$ that
meets only finitely many of the sets $X_\alpha$, we have
\[
  \dim_x(X)\leq\dim(U)=\sup_\alpha\bigl(\dim(U\cap X_\alpha)\bigr)
  \leq\sup_\alpha\bigl(\dim(X_\alpha)\bigr)
\]
by (14.1.4), whence the second assertion.
\end{proof}

\begin{corollary}[14.1.8]
\label{0.14.1.8}
Let $X$ be a Noetherian Kolmogorov space ($0_{\mathrm I}$, 2.1.3), and
let $F$ be the set of closed points of $X$. Then
$\dim(X)=\sup_{x\in F}\dim_x(X)$.
\end{corollary}

\begin{proof}
With the notation of the proof of (14.1.6), it suffices to remark
that there exists a closed point in $Z_0$ ($0_{\mathrm I}$, 2.1.3).
\end{proof}

\begin{proposition}[14.1.9]
\label{0.14.1.9}
Let $X$ be a non-empty Noetherian Kolmogorov space. For
$\dim(X)=0$, it is necessary and sufficient that $X$ be finite and discrete.
\end{proposition}

\begin{proof}
If a space $X$ is separated (and \emph{a fortiori} if $X$ is a discrete
space), the only irreducible closed sets of $X$ consist of a single
point, so $\dim(X)=0$. Conversely, suppose that $X$ is a Noetherian
Kolmogorov space such that $\dim(X)=0$; since every irreducible
component of $X$ contains a closed point ($0_{\mathrm I}$, 2.1.3), it
must reduce to that point. Since $X$ has only finitely many irreducible
components, it is finite and discrete.
\end{proof}

\begin{corollary}[14.1.10]
\label{0.14.1.10}
Let $X$ be a Noetherian Kolmogorov space. For a point $x\in X$ to be
isolated, it is necessary and sufficient that $\dim_x(X)=0$.
\end{corollary}

\begin{proof}
The condition is evidently necessary (without any hypothesis on $X$). It is
sufficient,
\oldpage[0\textsubscript{IV-1}]{8}
since it implies that $\dim(U)=0$ for an open neighbourhood $U$ of $x$, and
since $U$ is a Noetherian Kolmogorov space, $U$ is finite and discrete.
\end{proof}

\begin{proposition}[14.1.11]
\label{0.14.1.11}
The function $x\mapsto\dim_x(X)$ is upper semi-continuous on $X$.
\end{proposition}

\begin{proof}
It is clear that this function is upper semi-continuous at every
point where it takes the value $+\infty$. So suppose $\dim_x(X)=n<+\infty$; then
formula (14.1.4.1) shows that there exists an open neighbourhood $U_0$ of $x$ such
that $\dim(U)=n$ for every open neighbourhood $U\subset U_0$ of $x$. This being so,
for every $y\in U_0$ and every open neighbourhood $V\subset U_0$ of $y$, we have
$\dim(V)\leq\dim(U_0)=n$ (14.1.4); we therefore deduce from (14.1.4.1) that
$\dim_y(X)\leq n$.
\end{proof}

\begin{remark}[14.1.12]
\label{0.14.1.12}
If $X$, $Y$ are two topological spaces and $f:X\to Y$ is a continuous
map, we can have $\dim(f(X))>\dim(X)$; an example is obtained by taking
for $X$ a discrete space with 2 elements $a$, $b$, and for $Y$ the set
$\{a,b\}$ endowed with the topology for which the closed sets are
$\varnothing$, $\{a\}$ and $\{a,b\}$; if $f:X\to Y$ is the identity
map, we have $\dim(Y)=1$ and $\dim(X)=0$. Note that $Y$ is the spectrum
of a discrete valuation ring $A$, of which $a$ is the unique closed point and
$b$ the generic point; if $K$ and $k$ are the field of fractions and the
residue field of $A$, then $X$ is the spectrum of the ring $k\times K$ and $f$
is the continuous map corresponding to the homomorphism
$(\varphi,\psi):A\to k\times K$, where $\varphi:A\to k$ and $\psi:A\to K$ are
the canonical homomorphisms (cf. IV, 5.4.3).
\end{remark}

\subsection{Codimension of a closed subset}
\label{subsection:0.14.2}

\begin{definition}[14.2.1]
\label{0.14.2.1}
Given an irreducible closed subset $Y$ of a topological space $X$,
we call \emph{combinatorial codimension} (or simply \emph{codimension})
of $Y$ in $X$, and denote by $\operatorname{codim}(Y,X)$, the supremum
of the lengths of the chains of irreducible closed subsets of $X$ of which $Y$
is the smallest element. If $Y$ is an arbitrary closed subset of $X$, we
call \emph{codimension of $Y$ in $X$}, and again denote by
$\operatorname{codim}(Y,X)$, the infimum of the codimensions in $X$ of the
irreducible components of $Y$. We say that $X$ is
\emph{equicodimensional} if all the minimal irreducible closed subsets
of $X$ have the same codimension in $X$.
\end{definition}

It follows from this definition that
$\operatorname{codim}(\varnothing,X)=+\infty$, the infimum of the
empty subset of $\overline{\mathbb R}$ being $+\infty$. If $Y$ is closed in
$X$ and if $(X_\alpha)$ (resp. $(Y_\beta)$) is the family of irreducible
components of $X$ (resp. $Y$), every $Y_\beta$ is contained in some
$X_\alpha$, and more generally every chain of irreducible closed
subsets of $X$ of which $Y_\beta$ is the smallest element consists of
subsets of some $X_\alpha$; we therefore have
\[
  \operatorname{codim}(Y,X)=
  \inf_\beta\left(\sup_\alpha
  \bigl(\operatorname{codim}(Y_\beta,X_\alpha)\bigr)\right)
  \tag{14.2.1.1}\label{0.14.2.1.1}
\]
where for each $\beta$, $\alpha$ runs through the set of indices such that
$Y_\beta\subset X_\alpha$.

\begin{proposition}[14.2.2]
\label{0.14.2.2}
Let $X$ be a topological space.
\begin{enumerate}
  \item[\rm(i)] If $\Phi$ is the set of irreducible closed subsets of
  $X$, we have
  \[
    \dim(X)=\sup_{Y\in\Phi}\bigl(\operatorname{codim}(Y,X)\bigr).
    \tag{14.2.2.1}\label{0.14.2.2.1}
  \]
  \oldpage[0\textsubscript{IV-1}]{9}
  \item[\rm(ii)] For every non-empty closed subset $Y$ of $X$, we have
  \[
    \dim(Y)+\operatorname{codim}(Y,X)\leq\dim(X).
    \tag{14.2.2.2}\label{0.14.2.2.2}
  \]
  \item[\rm(iii)] If $Y$, $Z$, $T$ are three closed subsets of $X$ such
  that $Y\subset Z\subset T$, we have
  \[
    \operatorname{codim}(Y,Z)+\operatorname{codim}(Z,T)
    \leq\operatorname{codim}(Y,T).
    \tag{14.2.2.3}\label{0.14.2.2.3}
  \]
  \item[\rm(iv)] For a closed subset $Y$ of $X$ to satisfy
  $\operatorname{codim}(Y,X)=0$, it is necessary and sufficient that $Y$ contain an
  irreducible component of $X$.
\end{enumerate}
\end{proposition}

\begin{proof}
Assertions (i) and (iv) are immediate consequences of the definition
(14.2.1). To prove (ii), we may restrict to the case where $Y$ is
irreducible, and then the formula follows from the definitions (14.1.1) and
(14.2.1). Finally, to prove (iii), we may, by virtue of the definition
(14.2.1), first restrict to the case where $Y$ is irreducible; then
$\operatorname{codim}(Y,Z)=
\sup_\alpha(\operatorname{codim}(Y,Z_\alpha))$ for the irreducible
components $Z_\alpha$ of $Z$ containing $Y$; it is clear that
$\operatorname{codim}(Y,T)\geq\operatorname{codim}(Y,Z)$, so the inequality
is true if $\operatorname{codim}(Y,Z)=+\infty$; otherwise, there exists an
$\alpha$ such that
$\operatorname{codim}(Y,Z)=\operatorname{codim}(Y,Z_\alpha)$, and by virtue of
(14.2.1) we may restrict to the case where $Z$ too is irreducible; but
then the inequality (14.2.2.3) is an evident consequence of the definition
(14.2.1).
\end{proof}

\begin{proposition}[14.2.3]
\label{0.14.2.3}
Let $X$ be a topological space and $Y$ a closed subset of $X$. For every
open subset $U$ of $X$, we have
\[
  \operatorname{codim}(Y\cap U,U)\geq\operatorname{codim}(Y,X).
  \tag{14.2.3.1}\label{0.14.2.3.1}
\]
Moreover, for the two sides of (14.2.3.1) to be equal, it is necessary and
sufficient that, if $(Y_\alpha)$ is the family of irreducible components of
$Y$ meeting $U$, we have
$\operatorname{codim}(Y,X)=
\inf_\alpha(\operatorname{codim}(Y_\alpha,X))$.
\end{proposition}

\begin{proof}
We know ($0_{\mathrm I}$, 2.1.6) that $Z\mapsto\overline Z$ is a bijection
from the set of irreducible closed subsets of $U$ onto the set of
irreducible closed subsets of $X$ meeting $U$, and in particular maps
the irreducible components of $Y\cap U$ to the irreducible components
of $Y$ meeting $U$; if $Y_\alpha$ is one of the latter, we
therefore have
$\operatorname{codim}(Y_\alpha,X)=
\operatorname{codim}(Y_\alpha\cap U,U)$, and the proposition then follows from
the definition (14.2.1).
\end{proof}

\begin{definition}[14.2.4]
\label{0.14.2.4}
Let $X$ be a topological space, $Y$ a closed subset of $X$, and $x$ a point
of $X$. We call \emph{codimension of $Y$ in $X$ at the point $x$}, and denote by
$\operatorname{codim}_x(Y,X)$, the number
$\sup_U(\operatorname{codim}(Y\cap U,U))$, where $U$ runs through the set of
open neighbourhoods of $x$ in $X$.
\end{definition}

By virtue of (14.2.3), we can also write
\[
  \operatorname{codim}_x(Y,X)=
  \lim_U\bigl(\operatorname{codim}(Y\cap U,U)\bigr)
  \tag{14.2.4.1}\label{0.14.2.4.1}
\]
the limit being taken along the decreasing filtered set of open
neighbourhoods of $x$ in $X$. Note that we have
\[
  \operatorname{codim}_x(Y,X)=+\infty\quad\text{if}\quad x\in X-Y.
\]

\oldpage[0\textsubscript{IV-1}]{10}
\begin{proposition}[14.2.5]
\label{0.14.2.5}
If $(Y_i)_{1\leq i\leq n}$ is a finite family of closed subsets of a
topological space $X$, and $Y$ is the union of this family, we have
\[
  \operatorname{codim}(Y,X)=
  \inf_i\bigl(\operatorname{codim}(Y_i,X)\bigr).
  \tag{14.2.5.1}\label{0.14.2.5.1}
\]
\end{proposition}

\begin{proof}
Indeed, every irreducible component of one of the $Y_i$ is contained in
an irreducible component of $Y$, and conversely every irreducible
component of $Y$ is also an irreducible component of one of the $Y_i$
($0_{\mathrm I}$, 2.1.1); the conclusion therefore follows from the definition
(14.2.1) and from the inequality (14.2.2.3).
\end{proof}

\begin{corollary}[14.2.6]
\label{0.14.2.6}
Let $X$ be a topological space and $Y$ a locally Noetherian closed
subspace of $X$.
\begin{enumerate}
  \item[\rm(i)] For every $x\in X$, there exist only finitely many
  irreducible components $Y_i$ ($1\leq i\leq n$) of $Y$ containing $x$,
  and we have
  $\operatorname{codim}_x(Y,X)=
  \inf_i(\operatorname{codim}(Y_i,X))$.
  \item[\rm(ii)] The function
  $x\mapsto\operatorname{codim}_x(Y,X)$ is lower semi-continuous
  on $X$.
\end{enumerate}
\end{corollary}

\begin{proof}
Indeed, by hypothesis there is an open neighbourhood $U_0$ of $x$ in $X$ such
that $Y\cap U_0$ is Noetherian, and therefore has only finitely many irreducible
components, which are the traces on $U_0$ of irreducible components of $Y$;
\emph{a fortiori} there are only finitely many irreducible components
$Y_i$ ($1\leq i\leq n$) of $Y$ containing $x$, and, replacing $U_0$
by an open neighbourhood $U\subset U_0$ of $x$ meeting none of the $Y_j$
that do not contain $x$, we may suppose that the $Y_i\cap U$ are the irreducible
components of $Y\cap U$; for every open neighbourhood $V\subset U$ of $x$
in $X$, the $Y_i\cap V$ are then the irreducible components of
$Y\cap V$, and (14.2.3) then shows that
$\operatorname{codim}(Y_i,X)=\operatorname{codim}(Y_i\cap V,V)$, which
proves (i). Moreover, for every $x'\in U$, the irreducible components of
$Y$ containing $x'$ are among the $Y_i$, so
$\operatorname{codim}_{x'}(Y,X)\geq\operatorname{codim}_x(Y,X)$, which
proves (ii).
\end{proof}

\subsection{The chain condition}
\label{subsection:0.14.3}

\begin{env}[14.3.1]
\label{0.14.3.1}
In a topological space $X$, we will say that a chain
$Z_0\subset Z_1\subset\cdots\subset Z_n$ of irreducible closed subsets is
\emph{saturated} if there exists no irreducible closed subset $Z'$ distinct
from the $Z_i$ and such that $Z_k\subset Z'\subset Z_{k+1}$ for some index $k$.
\end{env}

\begin{proposition}[14.3.2]
\label{0.14.3.2}
Let $X$ be a topological space such that, for any two irreducible closed
subsets $Y$, $Z$ of $X$ with $Y\subset Z$, we have
$\operatorname{codim}(Y,Z)<+\infty$. The following two conditions are
equivalent:
\begin{enumerate}
  \item[\rm(a)] Two saturated chains of irreducible closed subsets of $X$
  having the same endpoints have the same length.
  \item[\rm(b)] If $Y$, $Z$, $T$ are three irreducible closed subsets of
  $X$ such that $Y\subset Z\subset T$, we have
  \[
    \operatorname{codim}(Y,T)=
    \operatorname{codim}(Y,Z)+\operatorname{codim}(Z,T).
    \tag{14.3.2.1}\label{0.14.3.2.1}
  \]
\end{enumerate}
\end{proposition}

\begin{proof}
It is immediate that a) implies b). Conversely, suppose that b) holds, and
let us prove that if two saturated chains with the same endpoints have lengths
$m$ and $n\leq m$, we necessarily have $m=n$. We argue by induction on
$n$, the proposition being evident for $n=1$. So suppose $n>1$, $n<m$,
and let $Z_0\subset Z_1\subset\cdots\subset Z_n$ be a
\oldpage[0\textsubscript{IV-1}]{11}
saturated chain such that there exists another saturated chain with endpoints
$Z_0$, $Z_n$ and of length $m$. Since
$\operatorname{codim}(Z_0,Z_n)\geq m>n$ and
$\operatorname{codim}(Z_0,Z_1)=1$, it follows from b) that
$\operatorname{codim}(Z_1,Z_n)=\operatorname{codim}(Z_0,Z_n)-1>n-1$, which
contradicts the induction hypothesis.
\end{proof}

When the conditions of (14.3.2) are satisfied, we say that $X$ satisfies
the \emph{chain condition}, or also that it is a \emph{catenary space}.
Every closed subspace of a catenary space is clearly catenary.

\begin{proposition}[14.3.3]
\label{0.14.3.3}
Let $X$ be a Noetherian Kolmogorov space of finite dimension. The
following conditions are equivalent:
\begin{enumerate}
  \item[\rm(a)] Any two maximal chains of irreducible closed subsets of $X$
  have the same length.
  \item[\rm(b)] $X$ is equidimensional, equicodimensional and catenary.
  \item[\rm(c)] $X$ is equidimensional, and for any two irreducible closed
  subsets $Y$, $Z$ of $X$ such that $Y\subset Z$, we have
  \[
    \dim(Z)=\dim(Y)+\operatorname{codim}(Y,Z).
    \tag{14.3.3.1}\label{0.14.3.3.1}
  \]
  \item[\rm(d)] $X$ is equicodimensional, and for any two irreducible closed
  subsets $Y$, $Z$ of $X$ such that $Y\subset Z$, we have
  \[
    \operatorname{codim}(Y,X)=
    \operatorname{codim}(Y,Z)+\operatorname{codim}(Z,X).
    \tag{14.3.3.2}\label{0.14.3.3.2}
  \]
\end{enumerate}
\end{proposition}

\begin{proof}
The hypotheses on $X$ imply that the endpoints of a maximal chain of
irreducible closed subsets of $X$ are necessarily a closed point and
an irreducible component of $X$ ($0_{\mathrm I}$, 2.1.3); moreover, every
saturated chain with endpoints $Y$, $Z$ (with $Y\subset Z$) is contained in a
maximal chain whose elements other than those of the given chain are
either contained in $Y$ or contain $Z$. These remarks
immediately establish the equivalence of a) and b), and also prove that if a)
holds, we have, for every irreducible closed subset $Y$ of $X$,
\[
  \dim(Y)+\operatorname{codim}(Y,X)=\dim(X)~;
  \tag{14.3.3.3}\label{0.14.3.3.3}
\]
since (14.3.2.1) holds, we immediately deduce (14.3.3.1) and (14.3.3.2) from
(14.3.3.3). Conversely, (14.3.3.1) implies (14.3.2.1), so (14.3.3.1)
implies the chain condition by virtue of (14.3.2); moreover, by
applying (14.3.3.1) to the case where $Y$ is reduced to a closed point $x$ of $X$ and
$Z$ to an irreducible component of $X$, we obtain
$\operatorname{codim}(\{x\},X)=\dim(Z)$; we conclude that c) implies b).
Likewise, (14.3.3.2) implies (14.3.2.1), hence the chain condition;
moreover, with the same choice of $Y$ and $Z$ as above, (14.3.3.2) again
implies $\operatorname{codim}(\{x\},X)=\dim(Z)$, so (every irreducible
component of $X$ containing a closed point by virtue of ($0_{\mathrm I}$,
2.1.3)) d) implies b).
\end{proof}

We say that a Noetherian Kolmogorov space is \emph{biequidimensional}
if it is of \emph{finite dimension} and satisfies the equivalent conditions
of (14.3.3).

\begin{corollary}[14.3.4]
\label{0.14.3.4}
Let $X$ be a biequidimensional Noetherian Kolmogorov space; then, for
every closed point $x$ of $X$ and every irreducible component $Z$ of $X$, we have
\[
  \dim(X)=\dim(Z)=\operatorname{codim}(\{x\},X)=\dim_x(X).
  \tag{14.3.4.1}\label{0.14.3.4.1}
\]
\end{corollary}

\oldpage[0\textsubscript{IV-1}]{12}
\begin{proof}
The last equality follows from the fact that if
$Y_0=\{x\}\subset Y_1\subset\cdots\subset Y_m$ is a maximal chain of
irreducible closed subsets of $X$ and $U$ is an open neighbourhood of $x$, the
$U\cap Y_i$ are pairwise distinct irreducible closed subsets of $U$
(since $\overline{U\cap Y_i}=Y_i$), whence
$\dim(U)=\dim(X)$ by virtue of (14.1.4).
\end{proof}

\begin{corollary}[14.3.5]
\label{0.14.3.5}
Let $X$ be a Noetherian Kolmogorov space; if $X$ is
biequidimensional, so is every union of irreducible
components of $X$ and every irreducible closed subset of $X$. Moreover,
for every closed subset $Y$ of $X$, we then have
\[
  \dim(Y)+\operatorname{codim}(Y,X)=\dim(X).
  \tag{14.3.5.1}\label{0.14.3.5.1}
\]
\end{corollary}

\begin{proof}
Since every chain of irreducible closed subsets of $X$ is contained in an
irreducible component of $X$, the first assertion follows immediately from
(14.3.3). Moreover, if $X'$ is an irreducible closed subset of $X$, $X'$
trivially satisfies the conditions (14.3.3, c)), whence the second assertion.

Finally, to prove (14.3.5.1), let us remark that we saw in the
proof of (14.3.3) that this relation holds when $Y$ is
irreducible; if $Y_i$ ($1\leq i\leq m$) are the irreducible components
of $Y$, the one among the $Y_i$ for which $\dim(Y_i)$ is largest is also
the one for which $\operatorname{codim}(Y_i,X)$ is smallest; therefore
(14.3.5.1) follows from the definitions of $\dim(Y)$ and of
$\operatorname{codim}(Y,X)$.
\end{proof}

\begin{remark}[14.3.6]
\label{0.14.3.6}
The reader will notice that the proof of (14.3.2) applies to an
arbitrary ordered set, the fact that we are dealing with irreducible closed
subsets of a topological space playing no role in that
proof. The same holds for the proof of (14.3.3) in an
ordered set $E$ such that for every $x\in E$ there exists a $z\leq x$ that
is \emph{minimal} in $E$, and in which the length of chains of elements of $E$
is bounded.
\end{remark}
